\section{模型假设}\label{sec:assumptions}

下列边界将题设条件与为求解所作的解释分开；模型性解释均以后文非前瞻回测和物理审计核验。“因果”只表示决策不使用决策时点后的信息，并非统计因果识别。

\subsection{共同物理边界}\label{sec:assumptions-common}
\begin{enumerate}[label=A\arabic*:,leftmargin=3.0em,labelsep=0.4em,itemsep=0.22em]
\item 题设设备与时间。一天含 $T=144$ 个 10 分钟时段，$\Delta t=1/6\,\mathrm{h}$；功率统一换算为能量 $L_t=\ell_t\Delta t$、$P_t=\rho_t\Delta t$。储能母线侧往返效率为 $90\%$，取 $\eta_c=\eta_d=\sqrt{0.9}$，$E_t\in[1200,10800]$ kWh，充、放电单时段能量均不超过 $5000\Delta t$ kWh。
\item 状态与能量流。2025-01-01 0:00 的 SOC 为 6000 kWh，跨日满足 $E_{d+1,0}=E_{d,T}$，不强制日末复位。购电、紧急购电、光伏与放电在公共母线供能；负荷、充电与弃光取能，普通购电额度不能直接跨时段挪用。
\item 适用范围。题目未给出售电价、PCC 容量、需量费、退化费和辅助功耗，主模型不引入这些量，也不另设购电功率上限；这限定了结论的工程边界。
\item 求解与审计。主模型为 LP。在正电价、无售电收益且有损耗时，同时充放电被更低损耗解支配；逐时段审计 $c_tr_t=0$，必要时以词典序次目标或 MILP 互斥复核。
\end{enumerate}

\subsection{各问信息、结算与输出边界}\label{sec:assumptions-q1}
\begin{enumerate}[label=Q\arabic*:,leftmargin=3.0em,labelsep=0.4em,itemsep=0.3em]
\item 日初一次性使用题给负荷、光伏预测制定全天计划，日内不更新；结果仅是确定性单日基准，不能与 Q2--Q4 的非前瞻策略混同。
\item 日初锁定 $q_{d,t}$，实际取用 $0\le x_{d,t}\le q_{d,t}$，并按 $p_tq_{d,t}+5p_te_{d,t}$ 结算，未用额度仍计费。预测和联合残差仅取目标日前已结束记录；每个 $(m,\alpha,K)$ 候选均以“情景计划—锁定 $q$—10 分钟因果 MPC—实际结算”的连续 SOC 闭环评分。$K=8$ 是该候选集下的经验证折中，而非全局最优情景数。1 月用于历史积累和 SOC 预热，题设工作簿及相应正文结果仅覆盖 2--12 月 334 日。
\item 0:00、6:00、12:00、18:00 接收预报，后三次更新从 6:10、12:10、18:10 生效且不改写已执行段；最终承诺 $g$ 只相对日初 $g^0$ 结算一次，$\phi_t=p_tg+\tfrac12p_t|g-g^0|$。日前负荷预测、光伏锚点插值和价值触发均只使用当时信息；终端价值不进入日账单，A 边界（年末 1200 kWh）为主口径，B 边界（6000 kWh）只作敏感性。
\item 实际 $p_{d,t}$ 是交付时段价格而非日前锁价：Q4-2 按 $p_{d,t}q_{d,t}+5p_{d,t}e_{d,t}$ 结算，Q4-3 的调整费也取交付时段价。价格预测训练集严格早于决策时点，联合情景沿用同一历史日的 $(r^L,r^P,r^p)$ 配对；仅在完成闭环校准、价格误差和物理/因果审计后采用正式结果。
\end{enumerate}
